\documentclass[a4paper,10pt]{article}
\usepackage[a4paper, margin=1in]{geometry}
\usepackage{amsmath,amssymb}
\usepackage{iftex}
\usepackage{enumitem}

\ifPDFTeX
  \usepackage[T1]{fontenc}
  \usepackage[utf8]{inputenc}
  \usepackage{textcomp}
  \usepackage{newtxtext,newtxmath}
\else
  \usepackage{fontspec}
  \setmainfont{Times New Roman}
\fi

\usepackage{lmodern}

\IfFileExists{upquote.sty}{\usepackage{upquote}}{}
\IfFileExists{microtype.sty}{%
  \usepackage[]{microtype}
  \UseMicrotypeSet[protrusion]{basicmath}
}{}

\makeatletter
\@ifundefined{KOMAClassName}{%
  \IfFileExists{parskip.sty}{%
    \usepackage{parskip}
  }{%
    \setlength{\parindent}{0pt}%
    \setlength{\parskip}{6pt}%
  }%
}{%
  \KOMAoptions{parskip=half}%
}
\makeatother

\usepackage{xcolor}
\usepackage{tabularx}
\usepackage{longtable}
\usepackage{booktabs}
\usepackage{titlesec}
\usepackage{bookmark}
\IfFileExists{xurl.sty}{\usepackage{xurl}}{}
\urlstyle{same}

\setlength{\emergencystretch}{3em}
\providecommand{\tightlist}{%
  \setlength{\itemsep}{0pt}\setlength{\parskip}{0pt}}
\setcounter{secnumdepth}{-\maxdimen}

\ifLuaTeX
  \usepackage{selnolig}
  \setmainfont{Times New Roman}
\fi

\hypersetup{
  hidelinks,
  pdfcreator={LaTeX via pandoc}
}

\author{}
\date{}

\titleformat{\subsection}
  {\normalfont\large\bfseries\color{black}}
  {\thesubsection}
  {1em}
  {}

\titleformat{\subsubsection}
  {\normalfont\normalsize\bfseries\color{black}}
  {\thesubsection}
  {1em}
  {}

\usepackage{fancyhdr}
\definecolor{darkgray}{gray}{0.4}
\fancypagestyle{firstpage}{
  \fancyhf{}
  \rhead{\textcolor{darkgray}{June 2026}}
  \renewcommand{\headrulewidth}{0pt}
}

\begin{document}
\thispagestyle{firstpage}

\begin{center}
\LARGE \textbf{\textsc{REUBEN H. KRAFT}} \\
\normalsize Professor of Mechanical and Biomedical Engineering, The Pennsylvania State University \\
\rule{\linewidth}{2pt}
\end{center}
\normalsize


\section*{HIGHLIGHTS OF RESEARCH ACCOMPLISHMENTS AND IMPACT}
\begin{itemize}[leftmargin=*, itemsep=0pt]
\item Research program impact includes \textbf{more than \$6.4 million in sponsored research funding, 69 peer-reviewed publications, mentorship of 20 graduate students (7 Ph.D. and 13 M.S.)}, and creation of sustained computational and experimental research capabilities spanning biomechanics, high strain-rate materials, protective systems, and digital engineering.
\item Recipient of the \textbf{Presidential Early Career Award for Scientists and Engineers (PECASE), the National Science Foundation CAREER Award, and election as Fellow of the American Society of Mechanical Engineers (ASME)} in recognition of contributions to computational and experimental mechanics.
\item Brings a unique perspective spanning federal laboratory, applied research, and academic environments through prior service as a \textbf{Mechanical Engineer at the U.S. Army Research Laboratory (ARL), Lead Researcher at the Johns Hopkins University Applied Physics Laboratory (APL), and Professor at The Pennsylvania State University}. This experience provides a rare understanding of the transition of fundamental research into operational and defense-relevant applications.
\item Founded \textbf{BrainSim Technologies, Inc.}, a Penn State spin-off company commercializing computational brain biomechanics technologies. Established a strategic partnership with \textbf{Wilson Sporting Goods, Inc.} to translate advanced brain injury simulation capabilities into consumer applications, with company revenues approaching \$1 million.
\item Established an internationally recognized research program in \textbf{computational and experimental mechanics of systems subjected to extreme dynamic loading}, integrating finite element simulation, high strain-rate experimentation, multiscale modeling, and digital engineering approaches to address challenges in survivability, protective systems, and human performance. Recognized as a trusted research leader by federal agencies and national laboratories through sustained sponsorship from the U.S. Army Research Laboratory, Air Force Research Laboratory, Los Alamos National Laboratory, Sandia National Laboratories, Department of Energy, National Science Foundation, National Institutes of Health, and multiple defense SBIR/STTR programs.
\item Coined and established the field of \textbf{Connectome Neurotrauma Mechanics}, integrating biomechanics, neuroscience, and brain connectomics to investigate how mechanical loading disrupts neural pathways and brain function. Developed pioneering \textbf{embedded finite element methodologies} for modeling axonal fiber tracts and traumatic brain injury under impact and blast loading.
\item Established and leads one of the nation's smallest \textbf{Split Hopkinson Pressure Bar (Kolsky Bar)} experimental facilities, enabling ultra-high strain-rate characterization of biological tissues and advanced engineering materials and providing critical validation data for predictive computational mechanics models.
\item Led a sustained research program with \textbf{Los Alamos National Laboratory} focused on deformation and failure of penetration-resistant composite materials, including Dyneema-based armor systems, integrating experimentation and simulation to improve protective technologies and survivability systems.
\item Developed predictive computational models and \textbf{digital-twin methodologies} for understanding long-term musculoskeletal degeneration and injury risk in military aviators, supporting Air Force efforts to improve pilot health, readiness, and performance.
\item Created \textbf{sensor-enabled and cloud-based simulation platforms} that integrate wearable sensors, experimental measurements, and physics-based computational models, enabling scalable digital engineering and digital twin capabilities for biomechanics and human-system applications.
\item Pursuing development of \textbf{hybrid quantum-computing approaches} for explicit dynamic finite element simulation, exploring next-generation computational architectures for large-scale predictive engineering analysis and future digital engineering environments.
\end{itemize}


\newpage


\section*{PROFESSIONAL POSITIONS}

\subsection*{Academic}

\noindent\begin{tabularx}{\linewidth}{@{}>{\raggedright\arraybackslash}X>{\raggedleft\arraybackslash}p{0.18\linewidth}@{}}Professor of Mechanical Engineering, The Pennsylvania State University, University Park, PA & 2024 - Present\\\end{tabularx}
\noindent\begin{tabularx}{\linewidth}{@{}>{\raggedright\arraybackslash}X>{\raggedleft\arraybackslash}p{0.18\linewidth}@{}}Associate Professor of Biomedical Engineering, The Pennsylvania State University University Park, PA & 2019 - 2024\\\end{tabularx}
\noindent\begin{tabularx}{\linewidth}{@{}>{\raggedright\arraybackslash}X>{\raggedleft\arraybackslash}p{0.18\linewidth}@{}}Associate Professor of Mechanical Engineering, The Pennsylvania State University University Park, PA & 2019 - 2024\\\end{tabularx}
\noindent\begin{tabularx}{\linewidth}{@{}>{\raggedright\arraybackslash}X>{\raggedleft\arraybackslash}p{0.18\linewidth}@{}}Assistant Professor of Biomedical Engineering, The Pennsylvania State University University Park, PA & 2016 - 2019\\\end{tabularx}
\noindent\begin{tabularx}{\linewidth}{@{}>{\raggedright\arraybackslash}X>{\raggedleft\arraybackslash}p{0.18\linewidth}@{}}Assistant Professor of Mechanical Engineering, The Pennsylvania State University University Park, PA & 2013 - 2019\\\end{tabularx}

\subsection*{Government}

\noindent\begin{tabularx}{\linewidth}{@{}>{\raggedright\arraybackslash}X>{\raggedleft\arraybackslash}p{0.18\linewidth}@{}}Mechanical Engineer, The U.S. Army Research Laboratory, Soldier Protection Sciences Branch & 2009 - 2012\\\end{tabularx}
\noindent\begin{tabularx}{\linewidth}{@{}>{\raggedright\arraybackslash}X>{\raggedleft\arraybackslash}p{0.18\linewidth}@{}}Post-Doc, Oak Ridge Associated Universities at The U.S. Army Research Laboratory, Impact Physics Branch & 2008 - 2009\\\end{tabularx}

\subsection*{Professional}

\noindent\begin{tabularx}{\linewidth}{@{}>{\raggedright\arraybackslash}X>{\raggedleft\arraybackslash}p{0.18\linewidth}@{}}Lead Researcher of Computational Biomechanics, The Johns Hopkins University Applied Physics Laboratory, Research and Exploratory Development Department, Biomechanics and Injury Mitigation Systems Group & 2012 - 2013\\\end{tabularx}




\section*{EDUCATION}

\noindent\begin{tabularx}{\linewidth}{@{}>{\raggedright\arraybackslash}X>{\raggedleft\arraybackslash}p{0.18\linewidth}@{}}Post-Doctoral, U.S. Army Research Laboratory, Aberdeen Proving Ground & 2008 - 2009\\\end{tabularx}
\noindent\begin{tabularx}{\linewidth}{@{}>{\raggedright\arraybackslash}X@{}}\textbf{Concentration:} Mechanics\\\end{tabularx}

\noindent\begin{tabularx}{\linewidth}{@{}>{\raggedright\arraybackslash}X>{\raggedleft\arraybackslash}p{0.18\linewidth}@{}}Ph D, The Johns Hopkins University, Baltimore & 2008\\\end{tabularx}
\noindent\begin{tabularx}{\linewidth}{@{}>{\raggedright\arraybackslash}X@{}}\textbf{Major:} Mechanical Engineering\\\end{tabularx}

\noindent\begin{tabularx}{\linewidth}{@{}>{\raggedright\arraybackslash}X>{\raggedleft\arraybackslash}p{0.18\linewidth}@{}}MS, The Johns Hopkins University, Baltimore & 2006\\\end{tabularx}
\noindent\begin{tabularx}{\linewidth}{@{}>{\raggedright\arraybackslash}X@{}}\textbf{Major:} Mechanical Engineering\\\end{tabularx}

\noindent\begin{tabularx}{\linewidth}{@{}>{\raggedright\arraybackslash}X>{\raggedleft\arraybackslash}p{0.18\linewidth}@{}}BS, University of Maryland, Baltimore County (UMBC) & 2003\\\end{tabularx}
\noindent\begin{tabularx}{\linewidth}{@{}>{\raggedright\arraybackslash}X@{}}\textbf{Major:} Mechanical Engineering\\\end{tabularx}




\section*{AWARDS AND HONORS}

\noindent\begin{tabularx}{\linewidth}{@{}>{\raggedright\arraybackslash}X>{\raggedleft\arraybackslash}p{0.18\linewidth}@{}}Professor of the Semester, American Society of Mechanical Engineers (ASME) Penn State Student Chapter & 2026\\\end{tabularx}
\noindent\begin{tabularx}{\linewidth}{@{}>{\raggedright\arraybackslash}X>{\raggedleft\arraybackslash}p{0.18\linewidth}@{}}PSEAS Premier Teaching Award, The Penn State Engineering Alumni Society (PSEAS) & 2025\\\end{tabularx}
\noindent\begin{tabularx}{\linewidth}{@{}>{\raggedright\arraybackslash}X>{\raggedleft\arraybackslash}p{0.18\linewidth}@{}}Outstanding Presentation Award, Penn State Center of Neural Engineering & 2023\\\end{tabularx}
\noindent\begin{tabularx}{\linewidth}{@{}>{\raggedright\arraybackslash}X>{\raggedleft\arraybackslash}p{0.18\linewidth}@{}}Fellow, American Society of Mechanical Engineering & 2023\\\end{tabularx}
\noindent\begin{tabularx}{\linewidth}{@{}>{\raggedright\arraybackslash}X>{\raggedleft\arraybackslash}p{0.18\linewidth}@{}}Best Poster Award, Penn State Center of Neural Engineering & 2022\\\end{tabularx}
\noindent\begin{tabularx}{\linewidth}{@{}>{\raggedright\arraybackslash}X>{\raggedleft\arraybackslash}p{0.18\linewidth}@{}}Web-Based Seminar, Inclusive Leadership in Equity, Allyship and Diversity (ILEAD) Training Series & 2022\\\end{tabularx}
\noindent\begin{tabularx}{\linewidth}{@{}>{\raggedright\arraybackslash}X>{\raggedleft\arraybackslash}p{0.18\linewidth}@{}}Faculty Early Career Development (CAREER) Program Award, National Science Foundation & 2019\\\end{tabularx}
\noindent\begin{tabularx}{\linewidth}{@{}>{\raggedright\arraybackslash}X>{\raggedleft\arraybackslash}p{0.18\linewidth}@{}}PSEAS Outstanding Teaching Award, The Penn State Engineering Alumni Society (PSEAS) & 2018\\\end{tabularx}
\noindent\begin{tabularx}{\linewidth}{@{}>{\raggedright\arraybackslash}X>{\raggedleft\arraybackslash}p{0.18\linewidth}@{}}First Place Paper and Oral Presentation (presented by my graduate student, I was senior author on paper), 13th Annual Penn State College of Engineering Research Symposium (CERS) & 2016\\\end{tabularx}
\noindent\begin{tabularx}{\linewidth}{@{}>{\raggedright\arraybackslash}X>{\raggedleft\arraybackslash}p{0.18\linewidth}@{}}People's Choice Poster Award (presented by my student, I was senior author on poster), 2016 (Ernst) Mach Conference & 2016\\\end{tabularx}
\noindent\begin{tabularx}{\linewidth}{@{}>{\raggedright\arraybackslash}X>{\raggedleft\arraybackslash}p{0.18\linewidth}@{}}Shuman Early Career Professorship, Penn State University Department of Mechanical and Nuclear Engineering & 2013\\\end{tabularx}
\noindent\begin{tabularx}{\linewidth}{@{}>{\raggedright\arraybackslash}X>{\raggedleft\arraybackslash}p{0.18\linewidth}@{}}Presidential Early Career Awards for Scientists and Engineers (PECASE), The White House; Office of Science and Technology Policy & 2011\\\end{tabularx}
\noindent\begin{tabularx}{\linewidth}{@{}>{\raggedright\arraybackslash}X>{\raggedleft\arraybackslash}p{0.18\linewidth}@{}}Professional Development Activities Related to Leadership and Service Responsibilities & \\\end{tabularx}




\section*{PUBLICATIONS}
\textit{Mentored student and postdoc co-authors are underlined.}

\subsection*{Journal Articles}
\begin{enumerate}
\item \underline{Reyes, A. N.}, DeWitt, T. R., \& \textbf{\textbf{Kraft,} R. H.} (2026). Personalization of the Toyota human model for safety (THUMS) using avatar-driven morphing for biomechanical simulations Biomechanics 6(37). Published. \url{https://doi.org/10.3390/biomechanics6020037}
\item Bagchi, A., Muci-Kuchler, K. H., \& \textbf{\textbf{Kraft,} R. H.} (2025). Editorial: Special issue on injury and damage mechanics for biological structures. Journal of Engineering and Science in Medical Diagnostics and Therapy 8(4). Published. \url{https://doi.org/10.1115/1.4069353}
\item \underline{Menghani, R.}, Bardall, C., Tanaka, M., \& \textbf{\textbf{Kraft,} R. H.} (2025). Exploring the connection between brain strain and cognitive changes: A Protocol Study. Journal of Engineering and Science in Medical Diagnostics and Therapy 8(4). Published. \url{https://doi.org/10.1115/1.4067765}
\item \underline{Hannah, T.}, \underline{Martin, V.}, Ellis, S., \& \textbf{\textbf{Kraft,} R. H.} (2024). High speed impact testing of UHMWPE composite using orthogonal arrays. Experimental Mechanics 64.  Front Cover Art Published. \url{https://doi.org/10.1007/s11340-024-01064-y}
\item \underline{Hannah, T.}, \underline{Martin, V.}, Ellis, S., \& \textbf{\textbf{Kraft,} R. H.} (2024). Influence of sampling rate on reproducibility and accuracy of miniature Kolsky bar experiments. Journal of Verification, Validation and Uncertainty Quantification 9(1). Published. \url{https://doi.org/10.1115/1.4065207}
\item \underline{Hannah, T.}, \underline{Martin, V.}, Ellis, S., \& \textbf{\textbf{Kraft,} R. H.} (2024). Impact of imperfect Kolsky bar experiments across different scales using finite elements. Journal of Verification, Validation and Uncertainty Quantification. Published. \url{https://doi.org/10.1115/1.4065206}
\item \underline{Martin, V.}, Hannah, T., Ellis, S., \& \textbf{\textbf{Kraft,} R. H.} (Corresponding Author) (2023). Using the embedded element finite element method to simulate impact of Dyneema plates. Fibers and Polymers. Published. \url{https://doi.org/10.1007/s12221-023-00417-z}
\item \underline{Hannah, T.}, Schuster, B., Baker, Z., Ellis, S., \& \textbf{\textbf{Kraft,} R. H.} Miniature Kolsky Bar Experiment Techniques Applied to UHMWPE Composite Analysis. Journal of Dynamic Behavior of Materials.
\item Zuidema, T. R., Bazarian, J. J., Kercher, K. A., Rettke, D. J., Mannix, R., \textbf{\textbf{Kraft,} R. H.}, Newman, S. D., Ejima, K., Steinfeldt, J. A., \& Kawata, K. (2023). Longitudinal association of clinical and biochemical biomarkers with head impact exposure in adolescent football. JAMA Network Open. Published. \url{https://doi.org/10.1001/jamanetworkopen.2023.16601}
\item \underline{Menghani, R. R.}, Dasans, A., \& \textbf{\textbf{Kraft,} R. H.} (Corresponding Author) (2023). A sensor-enabled cloud-based computing platform for computational brain biomechanics. Computer Methods in Biomechanics and Biomedical Engineering. Published. \url{https://doi.org/10.1016/j.cmpb.2023.107470}
\item Ramtani, S., Sanchez, J. F., Boucetta, A., \textbf{\textbf{Kraft,} R. H.}, Vaca-Gonzalez, J. J., \& Garzon-Alvarado, D. A. (2023). A coupled mathematical model between bone remodeling and tumors: a study of different scenarios using Komarova's model. Biomechanics and Modeling in Mechanobiology. Published. \url{https://doi.org/10.1007/s10237-023-01689-3}
\item Ji, S., Ghajari, M., Mao, H., \textbf{\textbf{Kraft,} R. H.}, Hajiaghamemar, M., Panzer, M. B., Willinger, R., Gilchrist, M. D., Kleiven, S., \& Stitzel, J. D. (2022). Use of brain biomechanical models for monitoring impact exposure in contact sports. Annals of Biomedical Engineering. Published. \url{https://doi.org/10.1007/s10439-022-02999-w}
\item \underline{Martin, V.}, \textbf{\textbf{Kraft,} R. H.} (Corresponding Author), \underline{Hannah, T.}, \& Ellis, S. (2022). An energy-based study of the embedded element method for explicit dynamics. Advanced Modeling and Simulation in Engineering Sciences. Published. \url{https://doi.org/10.1186/s40323-022-00223-x}
\item Adewole, D. O., Struzyna, L. A., Harris, J. P., Nemes, A. D., Burrell, J. C., Petrov, D., \textbf{\textbf{Kraft,} R. H.}, Chen, I., Serruya, M. D., Wolf, J. A., \& Cullen, K. (2021). Development of optically controlled "living electrodes" with long-projecting axon tracts for a synaptic brain-machine interface. Science Advances 7(4). Published. \url{https://doi.org/10.1126/sciadv.aay5347}
\item \underline{Marinov, T.}, \underline{Yuchi, L.}, Adewole, D. O., Cullen, D. Kacy, \& \textbf{\textbf{Kraft,} R. H.} (2020). A computational model of bidirectional axonal growth in micro-tissue engineered neuronal networks (micro-TENNs). In Silico Biology 13(3-4), pp. 85-99. Published. \url{https://doi.org/10.3233/ISB-180172}
\item \underline{Subramani, A. V.}, \underline{Whitley, P., Garimella, H. T.}, \& \textbf{\textbf{Kraft,} R. H.} (2020). Fatigue damage prediction in the annulus of cervical spine intervertebral discs using finite element analysis. Computer Methods in Biomechanics and Biomedical Engineering 23(11), 773-784. Published. \url{https://doi.org/10.1080/10255842.2020.1764545}
\item Carrera-Pinzon, A. F., Marquez-Florez, K., \textbf{\textbf{Kraft,} R. H.}, Ramtani, S., \& Garzon-Alvarado, D. A. (2019). Computational model of a synovial joint morphogenesis. Biomechanics and Modeling in Mechanobiology, 1--14. Published. \url{https://doi.org/10.1007/s10237-019-01277-4}
\item \underline{\textbf{\textbf{Kraft,} R. H.}}, \underline{Lee, C.}, Richtsmeier, J. T., \& \underline{Dolack, M. E.} (2019). Exploring mechanisms of cranial vault development using a coupled turing-biomechanical model. The FASEB Journal 33, 326.2-326.2. Published. \url{https://doi.org/10.1096/fasebj.2019.33.1_supplement.326.2}
\item \underline{Lee, C.}, Richtsmeier, J. T., \& \textbf{\textbf{Kraft,} R. H.} (2019). A coupled reaction-diffusion-strain model predicts cranial vault formation in development and disease. Biomechanics and Modeling in Mechanobiology. Published. \url{https://doi.org/10.1007/s10237-019-01139-z}
\item \underline{Przekwas, A. J., Tan, X. Gary, Chen, Z. J., Miao, Y., Harrand, V., Garimella, H. T.}, \textbf{\textbf{Kraft,} R. H.}, \& Gupta, R. K. (2019). Biomechanics of blast TBI with time resolved consecutive primary, secondary and tertiary loads. Military Medicine. Published. \url{https://doi.org/10.1093/milmed/usy344}
\item \underline{Garimella, H. T.}, \underline{Menghani, R.}, \underline{Gerber, J. I.}, \underline{Sridhar, S.}, \& \textbf{\textbf{Kraft,} R. H.} (2018). Embedded finite elements for modeling axonal injury. Annals of Biomedical Engineering. Published. \url{https://doi.org/10.1007/s10439-018-02166-0}
\item \underline{Motiwale, S.}, \underline{Subramani, A. V.}, Zhou, A., \& \textbf{\textbf{Kraft,} R. H.} (2018). A non-linear multi-axial fatigue damage model for the cervical intervertebral disc annulus. Advances in Mechanical Engineering 10(6). Published. \url{https://doi.org/10.1177/1687814018779494}
\item \underline{Dhobale, A. V.}, \underline{Adewole, O., Chan, A., Marinov, T.}, Serruya, M., \textbf{\textbf{Kraft,} R. H.}, \& Cullen, D. Kacy (2018). Assessing functional connectivity across 3D tissue engineered axonal tracts using calcium fluorescence imaging. Journal of Neural Engineering 15(5). Published. \url{https://doi.org/10.1088/1741-2552/aac96d}
\item \underline{Ranslow, A.}, \underline{Fang, Z.}, \underline{De Tomas, P.}, Gunnarsson, A., Weerasooriya, T., Satapathy, S., Thompson, K. A., \& \textbf{\textbf{Kraft,} R. H.} (2018). The multiaxial failure response of porcine trabecular skull bone estimated using microstructural simulations. American Society of Mechanical Engineers (ASME) Journal of Biomechanical Engineering 140(10). Published. \url{https://doi.org/10.1115/1.4039895}
\item \underline{Garimella, H. T.}, \textbf{\textbf{Kraft,} R. H.}, \& Przekwas, A. J. (2018). Do blast-induced skull flexures result in axonal deformation? PLOS One 13(3). Published. \url{https://doi.org/10.1371/journal.pone.0190881}
\item Serruya, M. D., Harris, J. P., Adewole, D. O., Struzyna, L. A., Burrell, J. C., Nemes, A., Petrov, D., \textbf{\textbf{Kraft,} R. H.}, Chen, H. I., Wolf, J. A., \& Cullen, D. K. (2017). Engineered axonal tracts as "living electrodes" for synaptic-based modulation of neural circuitry. Advanced Functional Materials, 1701183-n/a. Published. \url{https://doi.org/10.1002/adfm.201701183}
\item \underline{Lee, C. X.}, Richtsmeier, J. T., \& \textbf{\textbf{Kraft,} R. H.} (2017). A computational analysis of bone formation in the cranial vault using a coupled reaction-diffusion-strain model. Journal of Mechanics in Medicine and Biology 17(4). Published. \url{https://doi.org/10.1142/S0219519417500737}
\item \underline{Garimella, H. T.}, \& \textbf{\textbf{Kraft,} R. H.} (2017). A new computational approach for modeling diffusion tractography in the brain. Journal of Neural Regeneration Research 12(1). Published. \url{https://doi.org/10.4103/1673-5374.198967}
\item \underline{Garimella, H. T.}, \& \textbf{\textbf{Kraft,} R. H.} (2016). Modeling the mechanics of axonal fiber tracts using the embedded finite element method. International Journal for Numerical Methods in Biomedical Engineering 33(5), 1-21. Published. \url{https://doi.org/10.1002/cnm.2796}
\item \underline{Fielding, R. A.}, Przekwas, A. J., Tan, X. G., \& \textbf{\textbf{Kraft,} R. H.} (2015). Development of a lower extremity model for high strain rate impact loading. International Journal of Experimental and Computational Biomechanics 3(2), 161-186. Published. \url{https://doi.org/10.1504/IJECB.2015.070427}
\item \underline{Lee, C. X.}, Richtsmeier, J. T., \& \textbf{\textbf{Kraft,} R. H.} (2015). A computational analysis of bone formation in the cranial vault in the mouse. Frontiers in Bioengineering and Biotechnology 3(24). Published. \url{https://doi.org/10.3389/fbioe.2015.00024}
\item Swab, J. J., Tice, J., Wereszczak, A. A., \& \textbf{\textbf{Kraft,} R. H.} (2014). Fracture toughness of advanced structural ceramics: Applying ASTM C1421. Journal of the American Ceramic Society, pp. 1-9. Published. \url{https://doi.org/10.1111/jace.13293}
\item Clayton, J. D., \textbf{\textbf{Kraft,} R. H.}, \& Leavy, R. B. (2012). Mesoscale modeling of nonlinear elasticity and fracture in ceramic polycrystals under dynamic shear and compression. Journal of Solids and Structures 49(18), 6. Published. \url{https://doi.org/10.1016/j.ijsolstr.2012.05.035}
\item \textbf{\textbf{Kraft,} R. H.}, Mckee, P. J., Dagro, A. M., \& Grafton, S. T. (2012). Combining the finite element method with structural connectome-based analysis for modeling neurotrauma: Connectome neurotrauma mechanics. PLoS Computational Biology 8(8), e1002619. Published. \url{https://doi.org/10.1371%2Fjournal.pcbi.1002619}
\item \textbf{\textbf{Kraft,} R. H.}, \& Molinari, J. F. (2008). A statistical investigation of the effects of grain boundary properties on transgranular fracture. Acta Materialia 56(17), 10. Published. \url{https://doi.org/10.1016/j.actamat.2008.05.036}
\item \textbf{\textbf{Kraft,} R. H.}, Molinari, J. F., Ramesh, K. T., \& Warner, D. W. (2008). Computational micromechanics of dynamic compressive loading of a brittle polycrystalline material using a distribution of grain boundary properties. The Journal of Mechanics and Physics of Solids 56, 23. Published. \url{https://doi.org/10.1016/j.jmps.2008.03.009}
\end{enumerate}

\subsection*{Conference Proceedings}
\begin{enumerate}
\item \underline{Reyes Kadozono, A. N.}, Dewitt, T., \& \textbf{\textbf{Kraft,} R. H.} (2025). "Seat angle effects on disc degeneration for pilots in high-G environments." Proceedings of the 2024 American Society of Mechanical Engineers Congress and Exposition. 88629 .  Portland, Oregon, Nov. 17-20, 2024 Published. \url{https://doi.org/10.1115/IMECE2024-140657}
\item Leung, S. L., Greer, C. J., Byron, M., Ramos-Alvarado, B., \textbf{\textbf{Kraft,} R. H.}, Lesko, A. D., \& Berdanier, C. (2025). "Creating public resources to address resistance to diversifying content in mechanical engineering: Fostering awareness and ethical considerations." 2025 ASEE Annual Conference \& Exposition.  Montreal, Quebec, Canada, June 22-25 Nominated by the Mechanical Engineering Division for the 2025 ASEE Best Diversity, Equity \& Inclusion Paper Published. \url{https://doi.org/10.18260/1-2-56169}
\item \underline{Grube, R.}, Hannah, T. W., Martin, V. A., Ellis, S., \& \textbf{\textbf{Kraft,} R. H.} (2025). "Verification and validation of a novel miniaturized Kolsky bar using published Polymethyl Methacrylate (PMMA) data." Proceedings of the 2024 American Society of Mechanical Engineers Congress and Exposition. 88681 .  Portland, Oregon, Nov. 17-20, 2024 Published. \url{https://doi.org/10.1115/IMECE2024-145918}
\item Huber, C. M., Patton, D. A., Arbogast, K. B., \& \textbf{\textbf{Kraft,} R. H.} (2024). "Brain tissue strain during adolescent soccer heading using the cloud-based brain simulation research platform finite element head model." Proceedings of the 2024 International Research Council on Biomechanics on Injury. IRC-24-55 .
\item \underline{Hannah, T.}, \underline{\textbf{\textbf{Kraft,} R. H.}, Martin, V.}, \& Ellis, S. (2023). Impact of imperfect Kolsky bar experiments across different scales using finite elements.(IMECE2022-96816) . Proceedings of the 2022 American Society of Mechanical Engineers Congress and Exposition. Published. \url{https://doi.org/10.1115/IMECE2022-96816}
\item \underline{Martin, V.}, \underline{Hannah, T.}, Ellis, S., \& \textbf{\textbf{Kraft,} R. H.} (2023). Towards verification and validation of modeling Dyneema using the embedded finite element method.(IMECE2022-96784) . Proceedings of the 2022 American Society of Mechanical Engineers Congress and Exposition. Published. \url{https://doi.org/10.1115/IMECE2022-96784}
\item \underline{Hannah, T.}, \underline{\textbf{\textbf{Kraft,} R. H.}, Martin, V.}, \& Ellis, S. (2021). Implications of statistical spread to experimental analysis in a novel miniature Kolsky bar.(IMECE2020-23976) . Proceedings of the American Society of Mechanical Engineers Congress and Exposition.  Virtual, November 15-19, 2020 Published. \url{https://doi.org/10.1115/IMECE2020-23976}
\item \underline{Fang, Z.}, \underline{Ranslow, A. N.}, \& \textbf{\textbf{Kraft,} R. H.} (2016). Computational micromechanics of trabecular porcine skull bone using the material point method. Volume 3: Biomedical and Biotechnology Engineering(IMECE2016-67748), (pp. V003T04A044; 9 pages). Proceedings of the American Society of Mechanical Engineers Congress and Exposition.  Phoenix, Arizona, USA, November 11-17, 2016 Published. \url{https://doi.org/10.1115/IMECE2016-67748}
\item \underline{Motiwale, S.}, \underline{Subramani, V. V.}, Zhou, X., \& \textbf{\textbf{Kraft,} R. H.} (2016). Damage prediction for a cervical spine intervertebral disc. Volume 3: Biomedical and Biotechnology Engineering . Proceedings of the 2016 American Society of Mechanical Engineers Congress and Exposition.  Phoenix, Arizona, USA, November 11-17, 2016 Published. \url{https://doi.org/10.1115/IMECE2016-67711}
\item \underline{Chan, A. H. W., Dhobale, A.}, \underline{Adewole, O., Marinov, T.}, \underline{\textbf{\textbf{Kraft,} R. H.}}, Cullen, D. K., \& Serruya, M. (2016). Analysis of spontaneous calcium signals to infer functional connectivity within a novel "living electrode" neural construct. (pp. 1-2). Proceedings of IEEE.  Philadelphia, PA, USA, December 3, 2016. Published. \url{https://doi.org/10.1109/SPMB.2016.7846870}
\item \underline{Ranslow, A. N.}, \underline{\textbf{\textbf{Kraft,} R. H.}, Shannon, R.}, \underline{De Tomas-Medina, P.}, Radovitsky, R., Jean, A., Hautefeuille, M. P., Fagan, B., Ziegler, K. A., Weerasooriya, T., Dileonardi, A. M., Gunnarsson, A., \& Satapathy, S. (2016). Microstructural analysis of porcine skull bone subjected to impact loading. Volume 3: Biomedical and Biotechnology Engineering(IMECE2015-51979), (pp. V003T03A057; 10 pages). Proceedings of the American Society of Mechanical Engineers Congress and Exposition.  Houston, Texas, USA, November 13-19, 2015 Published. \url{https://doi.org/10.1115/IMECE2015-51979}
\item \underline{Lee, C.}, \& \textbf{\textbf{Kraft,} R. H.} (2016). A coupled reaction-diffusion-strain model for bone growth in the cranial vault. Proceedings of the 2016 Summer Biomechanics, Bioengineering and Biotransport Conference (SB3C2016). National Harbor, Maryland. June 29 - July 2, 2016
\item \underline{Ranslow, A. N.}, \& \textbf{\textbf{Kraft,} R. H.} (2016). The development of a "fuzzy" yield envelope for trabecular porcine skull bone using numerical simulations. Proceedings of the 2016 Summer Biomechanics, Bioengineering and Biotransport Conference (SB3C2016). National Harbor, Maryland. June 29 - July 2, 2016
\item \underline{Motiwale, S.}, Eppler, W., Hollingsworth, D., Hollingsworth, C., Morgenthau, J., \& \textbf{\textbf{Kraft,} R. H.} (2016). Application of neural networks for filtering non-impact transients recorded from biomechanical sensors. Proceedings of the Institute of Electrical and Electronic Engineers (IEEE) International Conference on Biomedical and Health Informatics. (pp. 204 - 207). IEEE.  Las Vegas, NV. February, 2016. Published. \url{https://doi.org/10.1109/BHI.2016.7455870}
\item \underline{Reddy, S. N.}, \underline{Fielding, R. A.}, \underline{Robinson, M. J.}, \& \textbf{\textbf{Kraft,} R. H.} (2015). A computational study of fracture in the calcaneus under variable impact conditions. Volume 3: Biomedical and Biotechnology Engineering(IMECE2015-51984), (pp. V003T03A058; 10 pages). Proceedings of the American Society of Mechanical Engineers Congress and Exposition.  Houston, Texas, USA, November 13-19, 2015 Published. \url{https://doi.org/10.1115/IMECE2015-51984}
\item \textbf{\textbf{Kraft,} R. H.}, \& \underline{Garimella, H. T.} (2015). Embedded finite elements for modeling traumatic axonal injury. Proceedings of the Summer Biomechanics, Bioengineering and Biotransport Conference (SB3C 2015). Snowbird, Utah, June 17-20, 2015
\item \underline{Fielding, R. A.}, \underline{Tan, X. G., Przekwas, A. J., Kozuch, C. D.}, \& \textbf{\textbf{Kraft,} R. H.} (2015). High rate impact to the human calcaneus: A micromechanical analysis. Volume 3: Biomedical and Biotechnology Engineering(IMECE2014-38930), (pp. V003T03A009, (8 pages)). Proceedings of the American Society of Mechanical Engineers Congress and Exposition.  Montreal, Canada, November 14 - 20, 2014. Published. \url{https://doi.org/10.1115/IMECE2014-38930}
\item \underline{Garimella, H. T.}, \underline{Yaun, H.}, Johnson, B. D., Slobounov, S., \& \textbf{\textbf{Kraft,} R. H.} (2014). Anisotropic constitutive model of human brain with intravoxel heterogeneity of fiber orientation using diffusion spectrum imaging (DSI). Volume 3: Biomedical and Biotechnology Engineering(IMECE2014-39107), (pp. V003T03A011; 9 pages). Proceedings of the 2014 American Society of Mechanical Engineers Congress and Exposition.  Montreal, Canada, November 14 - 20, 2014. Published. \url{https://doi.org/10.1115/IMECE2014-39107}
\item \underline{Makwana, A. R.}, \underline{Krishna, A. R.}, \underline{Yuan, H.}, \textbf{\textbf{Kraft,} R. H.}, Zhou, X., Przekwas, A. J., \& Whitley, P. (2014). Towards a micromechanical model of intervertebral disc degeneration under cyclic loading.(IMECE2014-39174), (pp. V003T03A012; 7 pages). Proceedings of the American Society of Mechanical Engineers Congress and Exposition.  Montreal, Canada, November 14 - 20, 2014. Published. \url{https://doi.org/10.1115/IMECE2014-39174}
\item \underline{Lee, C.}, Richtsmeier, J. T., \& \textbf{\textbf{Kraft,} R. H.} (2014). A multiscale computational model for the growth of the cranial vault in craniosynostosis.(IMECE2014-38728), (pp. V009T12A061; 6 pages). Proceedings of the American Society of Mechanical Engineers Congress and Exposition.  Montreal, Canada, November 14 - 20, 2014. Published. \url{https://doi.org/10.1115/IMECE2014-38728}
\item \underline{Fielding, R. A.}, \textbf{\textbf{Kraft,} R. H.}, Ryan, T. M., \& Stecko, T. D. (2014). A micromechanics-based simulation of calcaneus fracture and fragmentation due to impact loading. Proceedings of the 11th World Congress on Computational Mechanics (WCCM XI) 5th. European Conference on Computational Mechanics (ECCM V) 6th. European Conference on Computational Fluid Dynamics (ECFD VI).
\item Zhang, J., Merkle, A. C., Carneal, C. M., Armiger, R. S., \textbf{\textbf{Kraft,} R. H.}, Ward, E. E., Ott, K. A., Wickwire, A. C., Dooley, C. J., Harrigan, T. P., \& Roberts, J. C. (2013). Effects of torso-borne mass and loading severity on early response of the lumbar spine under high-rate vertical loading. International Research Council on Biomechanics of Injury. Gothenburg, Sweden, September 11 - 13, 2013
\item \textbf{\textbf{Kraft,} R. H.}, Dagro, A. M., McKee, P. J., Grafton, S. T., Vettel, J., McDowell, K., Vindiola, M., \& Merkle, A. C. (2013). Combining the finite element method with structural network-based analysis for modeling neurotrauma. (pp. 4). 11th International Symposium, Computer Methods in Biomechanics and Biomedical Engineering.
\item Scheidler, M., Fitzpatrick, J., \& \textbf{\textbf{Kraft,} R. H.} (2011). In Tom Proulx (Ed.), Optimal pulse shapes for SHPB tests on soft materials. 1, (pp. 259-268). Society for Experimental Mechanics Series, Dynamic Behavior of Materials.  ISBN/ISSN: 2191-5644 DOI 10.1007/978-1-4614-0216-9. June, Uncasville, Connecticut. Published. \url{https://doi.org/10.1007/978-1-4614-0216-9_37}
\item \textbf{\textbf{Kraft,} R. H.}, Lynch, M. L., \& Vogel, E. W. (2011). Computational failure modeling of lower extremities. RTO-MP-HFM-207AC/323(HFM-207)(TP/412) . NATO Human Factors and Medicine Panel. ISBN/ISSN: 978-92-837-0153-8
\item Clayton, J. D., \& \textbf{\textbf{Kraft,} R. H.} (2011). Mesoscale modeling of dynamic failure of ceramic polycrystals. Advances in Ceramic Armor VII: Ceramic Engineering and Science Proceedings. (32), (pp. 237-248). Proceedings of the 35th International Conference on Advanced Ceramics and Composites. ISBN/ISSN: 10.1002/9781118095256.ch21
\item Vettel, J. M., Bassett, D. S., \textbf{\textbf{Kraft,} R. H.}, \& Grafton, S. T. (2010). Physics-based models of brain structure connectivity informed by diffusion weighted imaging. Proceedings of the 27th Army Science Conference.
\item Gazonas, G. A., McCauley, J. W., \textbf{\textbf{Kraft,} R. H.}, Love, B. M., Clayton, J. D., Casem, D., Dandekar, D., Rice, B., Batyrev, I., Weingarten, N. S., \& Schuster, B. E. (2010). Multiscale modeling of armor ceramics: Focus on AlON. 27th Army Science Conference.
\item Scheidler, M., \& \textbf{\textbf{Kraft,} R. H.} (2010). In C. P. Hoppel (Ed.), Inertial effects in compression Hopkinson bar tests on soft materials. U.S. Army Research Laboratory, 1st Annual ARL Ballistic Technology Workshop.
\item \textbf{\textbf{Kraft,} R. H.}, Batyrev, I., Lee, S., Rollett, A. D., \& Rice, B. (2010). In J. J. Swab, S. Mathur and T. Ohji (Eds.), "Multiscale modeling of armor ceramics." Journal of the American Ceramics Society Meeting Proceedings. 31 . Hoboken, NJ: John Wiley \& Sons, Inc..  Ch.13. Published. \url{https://doi.org/10.1002/9780470944004}
\item Wereszczak, A. A., \& \textbf{\textbf{Kraft,} R. H.} (2003). In W. M. Kriven and H. T. Lin (Eds.), Flexural and torsional resonances of ceramic tiles via impulse excitation of vibration. 24(4), (pp. 207-213). 27th Annual Conference on Advanced Ceramics and Composites: B: Ceramic Engineering and Science Proceedings. Published. \url{https://doi.org/10.1002/9780470294826.ch31}
\item Wereszczak, A. A., \& \textbf{\textbf{Kraft,} R. H.} (2002). In H. T. Lin and M. Singh (Eds.), Instrumented Hertzian indentation of armor ceramics. 23(3), (pp. 11). 26th Annual Conference on Composites, Advanced Ceramics, Materials, and Structures: A: Ceramic Engineering and Science Proceedings. Published. \url{https://doi.org/10.1002/9780470294741.ch7}
\end{enumerate}

\subsection*{Preprints and Technical Reports}
\begin{enumerate}
\item \underline{Marinov, T.}, \underline{Yuchi, L.}, Adewole, D. O., Cullen, D. K., \& \textbf{\textbf{Kraft,} R. H.} "A computational model of bidirectional axonal growth in micro-tissue engineered neuronal networks (micro-TENNs)." bioRxiv. Cold Spring Harbor Laboratory. Published. \url{https://doi.org/10.1101/369843}
\item \underline{Gerber, J. I.}, \textbf{\textbf{Kraft,} R. H.}, \& \underline{Garimella, H. T.} (2018). "Computation of history-dependent mechanical damage of axonal fiber tracts in the brain: towards tracking sub-concussive and occupational damage to the brain." bioRxiv. Published. \url{https://doi.org/10.1101/346700}
\item \underline{Garimella, H. T.}, \underline{Menghani, R.}, \underline{Gerber, J. I.}, \underline{Sridhar, S.}, \& \textbf{\textbf{Kraft,} R. H.} (2018). "Embedded finite elements for modeling axonal injury." engrXiv. Published. \url{https://doi.org/10.31224/osf.io/2dx5e}
\item Adewole, D. O., Struzyna, L. A., Harris, J. P., Nemes, A. D., Burrell, J. C., Petrov, D., \textbf{\textbf{Kraft,} R. H.}, Chen, I., Serruya, M. D., Wolf, J. A., \& Cullen, K. (2018). "Optically-controlled "living electrodes" with long-projecting axon tracts for a synaptic brain-machine interface." bioRxiv. Published. \url{https://doi.org/10.1101/333526}
\item Dagro, A. M., McKee, P. J., \textbf{\textbf{Kraft,} R. H.}, Zhang, T. G., \& Satapathy, S. S. (2013). A preliminary investigation of traumatically induced axonal injury in a three-dimensional (3-D) finite element model (FEM) of the human head during blast-loading. Army Research Laboratory Technical Report (ARL-TR-6504).
\item Vettel, J., Dagro, A. M., Gordon, S., Kerick, S., \textbf{\textbf{Kraft,} R. H.}, Luo, S., Rawal, S., Vindiola, M., \& McDowell, K. (2012). Brain structure-function couplings (FY11). Army Research Laboratory Technical Report (ARL-TR-5893).
\item \textbf{\textbf{Kraft,} R. H.}, \& Wozniak, S. L. (2011). A review of computational spinal injury biomechanics research and recommendations for future efforts. Army Research Laboratory Technical Report (ARL-TR-5673).
\item \textbf{\textbf{Kraft,} R. H.}, \& Dagro, A. M. (2011). Design and implementation of a numerical technique to inform anisotropic hyperelastic finite element models using diffusion-weighted imaging. Army Research Laboratory Technical Report (ARL-TR-5796).
\item Clayton, J. D., \& \textbf{\textbf{Kraft,} R. H.} (2011). Mesoscale modeling of dynamic failure of ceramic polycrystals. Army Research Laboratory Reprint (ARL-RP-328).
\item Gozonas, G. A., McCauley, J. W., Batyrev, I. G., Casem, D., Clayton, J. D., Dandekar, D. P., \textbf{\textbf{Kraft,} R. H.}, Love, B. M., Rice, B. M., Schuster, B. E., \& Weingarten, N. S. (2011). Multiscale modeling of armor ceramics: Focus on AlON. Army Research Laboratory Reprint (ARL-RP-337).
\item Vettel, J. M., Bassett, D., \textbf{\textbf{Kraft,} R. H.}, \& Grafton, S. (2010). Physics-based models of brain structure connectivity informed by diffusion-weighted imaging. Army Research Laboratory Technical Reprint (ARL-RP-0355). Aberdeen Proving Ground, MD: U.S. Army Research Laboratory.
\item Wereszczak, A. A., Swab, J. J., \& \textbf{\textbf{Kraft,} R. H.} (2005). Effects of machining on the uniaxial and equibiaxial flexure strength of CAP3 AD-995 Al2O3. Army Research Laboratory Technical Report (ARL-TR-3617).
\item Swab, J. J., Wereszczak, A. A., Tice, J., Caspe, R., \textbf{\textbf{Kraft,} R. H.}, \& Adams, J. (2005). Mechanical and thermal properties of advanced ceramics for gun barrel applications. Army Research Laboratory Technical Report (ARL-TR-3417).
\end{enumerate}

\subsubsection*{Manuscripts Accepted for Publication}
\begin{enumerate}
\item \underline{Croom, J.}, Hannah, T., \& \textbf{\textbf{Kraft,} R. H.} "Direct optical measurement of radial strain in split Hopkinson pressure bar experiments." 2026 SEM Annual Conference and Exposition on Experimental and Applied Mechanics. [Accepted February 2026]. Norfolk, Virginia, June 1- 4, 2026
\end{enumerate}




\section*{CONTRACTS, FELLOWSHIPS, GRANTS, AND SPONSORED RESEARCH}

\noindent Kraft, Reuben H. (Principal Investigator), ``Elucidating High Strain Rate Deformation Mechanisms in Penetration-Resistant Composites'', Sponsored by Triad National Security, LLC (Los Alamos National Laboratory). June 2022 - September 2029.\vspace{0.25cm}

\noindent Kraft, Reuben H. (Principal Investigator), ``Quantifying, not exploring, the link between finite element-based strain predictions and cognitive changes'', Sponsored by Chuck Noll Foundation. January 2025 - February 2027.\vspace{0.25cm}

\noindent Kraft, Reuben H. (Principal Investigator), ``Development of Predictive Disc Degeneration Simulations for Pilots'', Sponsored by Air Force Research Laboratory. September 2015 - December 2025.\vspace{0.25cm}

\noindent Kraft, Reuben H. (Principal Investigator), ``Examining the link between finite element-based strain predictions and cognitive changes.'', Sponsored by Chuck Noll Foundation. August 2022 - August 2024.\vspace{0.25cm}

\noindent Kraft, Reuben H. (Principal Investigator), ``Unfunded Collaborative Research Agreement - Sports \& Wellbeing Analytics'', Sponsored by Sports \& Wellbeing Analytics. March 2021 - March 2024.\vspace{0.25cm}

\noindent Kraft, Reuben H. (Principal Investigator), ``CAREER: Multiscale Modeling of Axonal Fiber Bundles in the Brain'', Sponsored by National Science Foundation. February 2019 - January 2024.\vspace{0.25cm}

\noindent Kraft, Reuben H. (Co-Investigator), ``STTR PHASE II Enhancing Thermo-Mechanically Coupled Computational Models for High-Temperature Impact and Fracture'', Sponsored by Karagozian \& Case, Inc.. July 2021 - December 2023.\vspace{0.25cm}

\noindent Kraft, Reuben H. (Principal Investigator), ``Occupational mTBI from repeated exposure to low-level blast'', Sponsored by Defence Research and Development Canada (DRDC). May 2022 - March 2023.\vspace{0.25cm}

\noindent Kraft, Reuben H. (Principal Investigator), ``Development of a Novel Ballistic Armor Concept using FEM'', Sponsored by Triad National Security, LLC (Los Alamos National Laboratory). July 2017 - March 2022.\vspace{0.25cm}

\noindent Kraft, Reuben H. (Principal Investigator), ``An Exploration of the Material Point Method (MPM) in CTH Applied to Soft Material Systems Subjected to Dynamic Loading'', Sponsored by Sandia National Laboratories. January 2015 - December 2021.\vspace{0.25cm}

\noindent Kraft, Reuben H. (Principal Investigator), ``Head Kinematics Experimentation and Data Analysis'', Sponsored by U.S. Army Research Laboratory. September 2020 - May 2021.\vspace{0.25cm}

\noindent Kraft, Reuben H. (Co-Investigator), ``Craniosynostosis Network'', Sponsored by National Institutes of Health. February 2015 - January 2021.\vspace{0.25cm}

\noindent Kraft, Reuben H. (Co-Investigator), ``Stem Cell Mechanotransduction with Tendon Fatigue'', Sponsored by University of Pittsburgh. July 2018 - June 2020.\vspace{0.25cm}

\noindent Kraft, Reuben H. (Principal Investigator), ``SBIR Phase II: Global-Local Modeling of Aircraft Occupant Safety Assessment during Ejection (Air Force Phase II SBIR)'', Sponsored by CFD Research Corporation. October 2017 - January 2020.\vspace{0.25cm}

\noindent Kraft, Reuben H. (Principal Investigator), ``Development of Commercial Tools for Brain Modeling'', Sponsored by CFD Research Corporation. November 2017 - September 2019.\vspace{0.25cm}

\noindent Kraft, Reuben H. (Principal Investigator), ``NEUP: Multilayer Composite Fuel Cladding for LWR Performance Enhancement and Severe Accident Tolerance'', Sponsored by Department of Energy. October 2015 - June 2019.\vspace{0.25cm}

\noindent Kraft, Reuben H. (Principal Investigator), ``Embedded Finite Elements for a Multiscale, Multifunctional Approach for Modeling Axonal Bundles'', Sponsored by U.S. Army Research Laboratory. October 2016 - March 2019.\vspace{0.25cm}

\noindent Kraft, Reuben H. (Principal Investigator), ``Biological Living Electrodes Using Tissue Engineered Axonal Tracts to Probe and Modulate the Nervous System'', Sponsored by National Institutes of Health. September 2015 - July 2018.\vspace{0.25cm}

\noindent Kraft, Reuben H. (Principal Investigator), ``STTR Phase I: Synchronizing Video Imagery with Wearable Sensor Data and Side-by-Side Modeling Software to Develop Healthy Habits in Children'', Sponsored by CoachSafe PlaySafe, LLC. July 2016 - June 2018.\vspace{0.25cm}

\noindent Kraft, Reuben H. (Principal Investigator), ``Microstructural Analysis of Porcine Skull Bone Subjected to Impact Loading'', Sponsored by U.S. Army Research Laboratory. July 2014 - September 2017.\vspace{0.25cm}

\noindent Kraft, Reuben H. (Principal Investigator), ``SBIR: Phase II: A Neck Injury Assessment Tool for Prolonged Wear of Head Supported Mass'', Sponsored by CFD Research Corporation. January 2014 - June 2017.\vspace{0.25cm}

\noindent Kraft, Reuben H. (Principal Investigator), ``SBIR Phase II: Physics and Physiology Based Human Body Model of Blast Injury and Protection'', Sponsored by CFD Research Corporation. January 2014 - May 2017.\vspace{0.25cm}

\noindent Kraft, Reuben H. (Principal Investigator), ``Global-Local Modeling of Aircraft Occupant Safety Assessment during Ejection (Air Force SBIR Phase I)'', Sponsored by CFD Research Corporation. August 2016 - April 2017.\vspace{0.25cm}



\section*{INTELLECTUAL PROPERTY}

\noindent Kraft, R. H. "Brain Simulation Technology." Application Filed 2017.\newline
\noindent Kraft, R. H. "SmartGear: Instrumented Wrestling Headgear using Sensors." Application Filed 2017.\newline
\noindent Kraft, R. H. "Method and Apparatus for Multimodal Mobile Screening to Quantitatively Detect Brain Function Impairment," Application Filed September 2011.\newline




\section*{DIRECTED STUDENT LEARNING}
\subsection*{Ph.D. Dissertation}
\begin{enumerate}
\item Salas, Z, ``Computational high strain rate material characterization'', 2025 - Present
\item Croom, J, ``Experimental high strain rate material characterization'', 2024 - Present
\item Reyes, A, ``Modeling of spinal disc degeneration in fighter jet pilots'', 2022 - Present
\item Menghani, R, ``Sensor enabled, cloud-based modeling of the brain. Stage of Completion: Completed'', 2017 - May 2024
\item Martin, V, ``Modeling Armor Composites Undergoing High Strain Rate Deformation'', 2019 - August 2023
\item Hannah, T, ``Computational and experimental characterization of high strain rate response of Dyneema'', January 2018 - July 2023
\item Hertel, Z, ``An exploration of the Material Point Method (MPM) in CTH applied to soft material systems subjected to dynamic loading'', January 2015 - April 2023
\item Subramani, V, ``Modeling of spinal injury under extreme loading Conditions with emphasis on military loading Scenarios - a mathematical fatigue damage model and finite element study'', November 2015 - August 2020
\item Lee, C, ``A computational analysis of bone formation in the cranial vault using a reaction-diffusion-strain model'', December 2013 - May 2018
\item Garimella, H, ``An embedded element based human head model to investigate axonal injury'', September 2013 - June 2017
\end{enumerate}

\subsection*{Master's Thesis}
\begin{enumerate}
\item Paulin, D, ``Finite element modeling instrumented mouthguards'', 2024 - 2025
\item Grube, R, ``High strain rate material response of Dyneema'', 2023 - 2025
\item Fournier, N, ``Finite element modeling of gasket interfaces. Stage of Completion: Completed'', November 2023 - May 2024
\item Lovett, J, ``Energy based body armor design'', August 2021 - August 2023
\item Norris, I, ``Computational modeling of spinal degeneration in F35 pilots'', November 2020 - December 2022
\item Dolack, M, ``Computational morphogenesis of embryonic bone development: past, present, and future'', September 2017 - May 2019
\item Gerber, J, ``Development of a history-dependent damage model for the brain due to repetitive impacts'', November 2016 - May 2018
\item Dhobale, A, ``Assessing functional connectivity of micro-tissue engineered neural networks using calcium fluorescence imaging'', August 2016 - May 2017
\item Yuchi, L, ``A computational model of bidirectional growth for micro-Tissue Engineered Neuronal Networks (micro-TENNs)'', August 2016 - May 2017
\item Fang, Z, ``MPM methods for modeling trabecular bone'', August 2016 - May 2017
\item Motiwale, S, ``Modeling intervertebral disc degeneration due to cyclic loading'', January 2015 - May 2016
\item Ranslow, A, ``Microstructural analysis of porcine skull bone subjected to impact loading'', July 2014 - May 2016
\item Fielding, R, ``Development of a lower extremity model for high strain rate impact loading'', September 2013 - May 2015
\end{enumerate}

\subsection*{Postdoctoral Mentorship}
\begin{enumerate}
\item Marinov, T. Computational neuroscience: simulation of micro-tissue engineered neural networks., September 2016 - July 2018
\end{enumerate}




\section*{SERVICE}
\subsection*{College}

\noindent Engineering laptop Initiative, Member. July 2021 - December 2021. Developed laptop recommendations for college.\newline
\noindent Activity Insight Faculty Users Committee, Member. October 2017 - December 2020.\newline

\subsection*{Department}

\noindent Mechanical Engineering Strategic Plan Tracking Committee, Chairperson. August 2023 - May 2025.\newline
\noindent Promotion and Tenure Committee, Member. August 2023 - May 2024.\newline
\noindent Research Advancement Committee, Member. August 2023 - May 2024.\newline
\noindent Teaching Load Policy Committee, Chairperson. September 2022 - May 2023.\newline
\noindent Department Facilities Committee, Member. August 2022 - May 2023.\newline
\noindent Promotion and Tenure Committee, Member. August 2022 - May 2023.\newline
\noindent Research Advancement Committee, Member. August 2022 - May 2023.\newline
\noindent Mechanical Engineering Promotion and Tenure Committee, Committee Member. July 2022 - Present.\newline
\noindent Mechanical Engineering Strategic Planning Committee, Committee Member. January 2022 - December 2022.\newline
\noindent Teaching Load Policy Committee, Chairperson. September 2020 - May 2021.\newline
\noindent Mechanical Engineering Strategic Planning Committee, Member. August 2019 - September 2020.\newline
\noindent Joint Faculty Search in Mechanical Engineering and the Institute for CyberScience, Chairperson. August 2018 - May 2019.\newline
\noindent Faculty Search Committee for Mechanical Systems in Mechanical Engineering, Member. August 2017 - 2018.\newline
\noindent Mechanical Engineering Liaison to Institute for CyberScience, Liaison. 2017 - 2018.\newline
\noindent Faculty Search Committee for Emerging Areas in Mechanical Engineering, Member. August 2016 - 2017.\newline
\noindent Faculty Search Committee for Mechanical Systems in Mechanical Engineering, Member. September 2014 - March 2015.\newline

\subsection*{University}

\noindent Institute for Computational \& Data Sciences Coordinating Committee, Chairperson. August 2023 - May 2024.\newline
\noindent Graduate Council Committee on Academic Standards, Committee Member. June 2023 - Present.\newline
\noindent Graduate Council Representative to Engineering Faculty, Committee Member. June 2023 - Present.\newline
\noindent Institute for Computational \& Data Sciences Coordinating Committee, Co-Chairperson. August 2022 - May 2023.\newline
\noindent Hiring Committee for Project Coordinator for Institute for CyberScience, Member. March 2019 - April 2019.\newline




\section*{EDITORIAL BOARD POSITIONS}

\noindent ASME Journal of Engineering and Science in Medical Diagnostics and Therapy (JESMDT), Guest Editor. November 2024 - December 2025. Special issue on injury and damage mechanics for biological structures.\newline
\noindent ASME Journal of Engineering and Science in Medical Diagnostics and Therapy (JESMDT), Associate Editor. November 2022 - Present.\newline
\noindent Frontiers in Bioengineering and Biotechnology, Associate Editor. November 2014 - Present.\newline




\section*{PROFESSIONAL MEMBERSHIPS}

\noindent Member, United States Association for Computational Mechanics. February 2014 - 2015.\newline
\noindent Member, American Physical Society. May 2013 - 2014.\newline
\noindent Member, American Society for Engineering Education. May 2013 - 2014.\newline
\noindent Member, American Society of Biomechanics. May 2013 - 2014.\newline
\noindent Member, American Society of Mechanical Engineers. January 2003 - Present.\newline



\end{document}
